\documentclass[
    11pt,
    a4paper,
    oneside
]{article}

% ---------- Language and typography ----------
\usepackage[T1]{fontenc}
\usepackage[utf8]{inputenc} % Omit when using LuaLaTeX/XeLaTeX
\usepackage[english]{babel}
\usepackage{lmodern}
\usepackage{microtype}

% ---------- Page layout ----------
\usepackage[
    top=25mm,
    bottom=25mm,
    left=30mm,
    right=25mm
]{geometry}
\usepackage{setspace}
\onehalfspacing
\usepackage{parskip}
\setlength{\parindent}{0pt}

% ---------- Mathematics ----------
\usepackage{amsmath,amssymb,mathtools}
\usepackage{bm}
\usepackage{siunitx}
\sisetup{
    separate-uncertainty=true,
    per-mode=symbol,
    range-phrase=--,
    range-units=single
}

% Useful mathematical operators
\DeclareMathOperator{\diag}{diag}
\DeclareMathOperator{\RePart}{Re}
\DeclareMathOperator{\ImPart}{Im}
\newcommand{\dd}{\mathop{}\!\mathrm{d}}
\newcommand{\e}{\mathrm{e}}
\newcommand{\iu}{\mathrm{i}}
\newcommand{\vect}[1]{\bm{#1}}
\usepackage{amsmath,amssymb,mathtools}


\usepackage{amsmath}
\usepackage{graphicx}
\usepackage[europeanresistors]{circuitikz}
\ctikzset{bipoles/length=1.3cm}

% Circuit notation
\newcommand{\seriesconn}{\mathbin{\oplus}}
\newcommand{\Lel}[1]{\mathrm{L}\!\left(#1\right)}
\newcommand{\Wel}[1]{\mathrm{W}\!\left(#1\right)}
\newcommand{\Gel}[2]{\mathrm{G}\!\left(#1,#2\right)}
\newcommand{\CPEel}[2]{\mathrm{CPE}\!\left(#1,#2\right)}


% Grammar notation
\newcommand{\nt}[1]{\mathit{#1}}
\newcommand{\term}[1]{\text{\texttt{#1}}}
\newcommand{\defines}{\Coloneqq}
% Fractional derivatives
\newcommand{\Caputo}[2]{%
    {}^{\mathrm{C}}\!D_{#1}^{#2}%
}

% ---------- Figures and tables ----------
\usepackage{graphicx}
\graphicspath{{figures/}}
\newcommand{\reportgraphic}[2][]{%
    \IfFileExists{#2}{%
        \includegraphics[#1]{#2}%
    }{%
        \IfFileExists{figures/#2}{%
            \includegraphics[#1]{figures/#2}%
        }{%
            \PackageWarning{report}{Image `#2' not found; using a placeholder}%
            \fbox{%
                \parbox[c][45mm][c]{0.9\linewidth}{%
                    \centering Missing image: \texttt{\detokenize{#2}}%
                }%
            }%
        }%
    }%
}
\usepackage{booktabs}
\usepackage{tabularx}
\usepackage{multirow}
\usepackage{array}
\usepackage{subcaption}
\usepackage{float}

% ---------- Algorithms and source code ----------
\usepackage{algorithm}
\usepackage{algpseudocode}
\usepackage{listings}
\usepackage{xcolor}

\definecolor{codegreen}{RGB}{0,120,70}
\definecolor{codegray}{RGB}{100,100,100}
\definecolor{codeblue}{RGB}{20,70,160}
\definecolor{codebackground}{RGB}{247,247,247}

\lstset{
    backgroundcolor=\color{codebackground},
    basicstyle=\ttfamily\small,
    keywordstyle=\color{codeblue}\bfseries,
    commentstyle=\color{codegreen},
    stringstyle=\color{red!70!black},
    numbers=left,
    numberstyle=\tiny\color{codegray},
    frame=single,
    breaklines=true,
    showstringspaces=false,
    tabsize=4,
    captionpos=b
}

% ---------- References ----------
\usepackage[
    backend=biber,
    style=numeric-comp,
    sorting=none,
    giveninits=true,
    doi=true,
    url=true
]{biblatex}
\addbibresource{references.bib}

% ---------- Cross-references and links ----------
\usepackage{csquotes}
\usepackage{hyperref}
\hypersetup{
    colorlinks=true,
    linkcolor=blue!50!black,
    citecolor=green!40!black,
    urlcolor=blue!60!black,
    pdftitle={Fractional-Order za SOFCs},
    pdfauthor={Your Name}
}
\usepackage[nameinlink,capitalise,noabbrev]{cleveref}

% ---------- Headers and footers ----------
\usepackage{fancyhdr}
\pagestyle{fancy}
\fancyhf{}
\fancyhead[L]{FODE equation discovery}
\fancyhead[R]{\nouppercase{\leftmark}}
\fancyfoot[C]{\thepage}
\setlength{\headheight}{14pt}

% ---------- Section formatting ----------
\usepackage{titlesec}
\titleformat{\chapter}[hang]
    {\normalfont\huge\bfseries}
    {\thechapter.}
    {0.6em}
    {}

% ---------- Document metadata ----------
\title{
    \textbf{Fractional-Order Differential Equation Modelling\\
    of Solid Oxide Fuel Cells}
}
\author{
    Your Name\\
    \small Institution or Department
}
\date{\today}

\begin{document}

\section{Introduction}

In this report we present progress of formal grammar based equation discovery applied to EIS data in frequency domain,
both from solid oxide fuel cells and batteries.

%\section{Uvod}
%
%Po prejšnjem poročilu in okvirnem potrdilu, da se premikam v pravo smer, sem se lotil izboljšav algoritma. Dodal sem normalizacijo residualov (na vsakem koraku jih 
%delim s fiksno vrednostjo), in pa implementacija testiranja več začetnih vrednosti za vsako vezje. Rezultati so vzpodbudni, predvsem so se izboljšali pri vezjih brez 
%omejitev. %Znova bom razdelil poročilo na dva dela glede na dve gramatiki, ki jih testiram.

\section{The grammar}

We use formal grammars to define candidate circuits which we then fit on the data. We use a grammar that does not impose any structural
assumptions on the circuit. Each element can appear at any position, the only restriction being that a resistor $\term{Rs}$ is in series
with the rest of the circuit. The circuit elements present are standard equivalent circuit modeling (ECM) elements, which we present below.

%\section{Splošna gramatika}
%
%Osredotočil sem se na gramatiko brez stukturnih omejitev pri generaciji vezij, omejuje jo le število elementov, ki lahko nastopajo v vezju. Edina
%predpostavka je, da vsako vezje vsebuje en zaporedno vezan upornik $\term{Rs}$.

%Na grafu spodaj je predstavljen rezultat 5 najbolje prilegajočih vezij po BIC kriteriju izmed vseh 4356 vezij z največ 6 elementi. 

\begin{equation}
\label{eq:grammar-relaxed}
\begin{aligned}
\nt{Circuit}
    &\longrightarrow
    \term{Rs}
    \mid
    \term{Rs}\;\term{+}\;\nt{Network},
    \\[2mm]
\nt{Network}
    &\longrightarrow
    \nt{Element}
    \\%\mid
    &\longrightarrow\;\nt{Network}\;\term{+}\;\nt{Network}\;
    %&\longrightarrow\term{(}\;\nt{Network}\;\term{+}\;\nt{Network}\;\term{)}
    \\
    &\longrightarrow%\quad%\mid
    \;\nt{Network}\;\term{||}\;\nt{Network}\;,
    %\term{(}\;\nt{Network}\;\term{||}\;\nt{Network}\;\term{)},
    \\[2mm]
\nt{Element}
    &\longrightarrow
    \term{R}
    \mid \term{L}
    \mid \term{CPE}
    \mid \term{W}.
\end{aligned}
\end{equation}

%Impedance can be modelled by connecting series and parallel branches of basic elements. 
%\subsubsection*{Constant phase element (CPE)}
%Constant phase element is a component that models an imperfect capacitor. 
%\subsubsection*{Inductor (L) and Resistor (R)}
%
%Inductors and resistors 
%
%\subsubsection*{Warburg element (W)}


\section{Methodology}

Several steps were taken in circuit generation, subsequent derivation of impedance equations and finally parameter fitting.
Here we describe the core procedure.

\subsection{Circuit generation}

Circuit generation from formal grammars feels almost like solving a puzzle backwards. We consequently connect elements in series or parallel
and in the end map each $\nt{Element}$ to either a resistor, inductor, warburg element or a CPE. However, we can take a few measures to 
keep the number of circuit parameters as small as possible, namely:
\begin{itemize}
    \item replacing two resistors in series or in parallel with a single resistor
    \item replacing two inductors in series or in parallel with a single inductor
\end{itemize}

\subsection{Impedance equation generation}

Obtaining impedance equations from circuits generated by the grammar is straightforward. For series of elements we concatenate their
respective impedances, and for elements in parallel we sum their admittances. Below we list the elements present in the grammar and their impedances.

\subsubsection*{Constant phase element (CPE)}
Constant phase element is a component that models an imperfect capacitor. It contains a dimensionless parameter $0\le\alpha\leq1$ that can be
interpreted as a measure of non-ideality. When $\alpha=1$ a CPE is a perfect capacitor.
CPE's impedance is
\[
    Z_{CPE} = \frac{1}{(j\omega)^{\alpha}Q}
\]
where $j$ is the imaginary unit, and $Q$ is a parameter independent of the frequency. $\omega=2\pi f$ is the angular frequency
\subsubsection*{Inductor (L) and Resistor (R)}

Resistor's impedance is constant does not depend on the frequency, and inductor's impedance is a linear function of the frequency. 
$$Z_R = R,$$
where $R$ is the resistor's resistance, and
\[
    Z_L = j\omega L,
\]
with $L$ the inductance.
\subsubsection*{Warburg element (W)}
Warburg element is used mainly in batteries EIS analysis. It is a special case of the CPE where $\alpha=0.5$. We focused on the infinite Warburg
element with impedance defined as
\[
    Z_W = \frac{\sigma}{\sqrt{jw}}
\]
\subsection{Parameter fitting}
As described above, frequency domain impedance equations are directly generated from the circuits by simply summing up elements
impedance formulas when they are connected in series, and summing up their reciprocal values when they are connected in parallel circuits.

We set a ceiling of the number of elements that may feature in the circuits. We run the optimization algorithm (more on that in a later section)
up to a certain threshold, and for the 5 best fitting circuits based on mean squared error we re-run the optimization algorithm with 
stricter thresholds to obtain results as close to full convergence as possible. We sort these 5 based on standard BIC to include some 
penalty for circuit complexity.

\subsubsection{Initial value selection}

We opted for a multi start stratified approach. Each circuit gets fitted 6 times in the initial evaluation. First fit is using deterministic
fixed values $p_{fixed}$:

\begin{itemize}
    \item 0.8 for $\alpha$
    \item 0.05 for $R_s$
    \item 0.02 for other resistors $R$
    \item 1.0 for $Q$
    \item 0.01 for $\sigma$
    \item $10^{-6}$ for $L$
\end{itemize}

For the remaining 5 starts we selected the values as follows:
\begin{itemize}
    \item $\alpha$: Divide interval $[0.03,0.98]$ in 5 equal parts and sample from each of the subintervals
    \item $R_s$: Divide interval $[0.05\cdot10^{-1},0.05\cdot10]=[0.005,0.5]$ into 5 parts and sample from subintervals
    \item For the other parameters: Divide $[p_{fixed}\cdot10^{-2},p_{fixed}\cdot10^{2}]$ in 5 subintervals and sample from them
\end{itemize}

This way we get a decent spread of initial values. We chose a narrower band for $R_s$ as it is more directly constrained by the higher frequencies,
and kept a broader band to allow for more uncertainty in initial values.

\subsubsection{Numerical algorithm}\label{numalg}

In order to improve results of the algorithm we tried to improve robustness of SciPy's \texttt{least\_squares} algorithm. 
For the initial optimization run we use:

\begin{table}[]
\centering
\begin{tabular}{l|l}
Parameter & value \\ \hline
\texttt{x\_scale} & \texttt{jac} \\
\texttt{ftol} &  10{-6}\\
\texttt{xtol} &  10{-6}\\
\texttt{gtol} &  10{-6}\\
\texttt{f\_scale} &  0.001\\
\texttt{max\_nfev} &  100\\
\texttt{loss} &  \texttt{soft\_l1}\\
\texttt{method} & \texttt{trf}
\end{tabular}
\label{parsopt1}
\caption{Parameters for the SciPy.optimize least\_squares algorithm in the first optimization}
\end{table}

TRF method means we use the \textit{Trust region field} method of nonlinear least squares, and JAC scale is a specific 
way of calculating matrix column scales. Tolerance fields (f-,x-,g-) stop the iteration if function values, parameters or function's derivatives
move less than the threshold. Soft L1 loss with corresponding f-scale parameter means that instead of minimizing squared residuals, the optimizer
is instead minimizing the following objective function:

\begin{equation}
    F(\boldsymbol{x})
    = \frac{C^2}{2}
      \sum_{i=1}^{m}
      \rho\left(\frac{r_i(\boldsymbol{x})^2}{C^2}\right),
\end{equation}
where $r_i(\boldsymbol{x})$ denotes the $i$th residual.
For the \texttt{soft\_l1} loss with \texttt{f\_scale = 0.001},
\begin{equation}
    \rho(z) = 2\left(\sqrt{1+z}-1\right),
    \qquad C = 10^{-3}.
\end{equation}

Simply put, more residuals behave quadratically with a larger f-scale, while a smaller f-scale means most of the function is behaving like
the L1 loss. For the first optimization round, so for all circuits generated, the stopping point was the 100 function evaluations.

Substituting these expressions gives the optimization problem
\begin{equation}
    \boxed{
        \min_{\boldsymbol{x}} F(\boldsymbol{x})
        =
        \min_{\boldsymbol{x}}
        10^{-6}\sum_{i=1}^{m}
        \left[
            \sqrt{1+10^{6}r_i(\boldsymbol{x})^2}-1
        \right].
    }
\end{equation}

In the second optimization round, that is, for 5 circuits with smallest MSE values, we again used the same algorithm, but with stricter
optimization values presented in Table \ref{parsopt2}. Main difference is the increased number of steps and stricter function evaluation.
In this case, for each of the circuits one of the tolerances was reached before we hit the step limit.
\begin{table}[]
\centering
\begin{tabular}{l|l}
Parameter & value \\ \hline
\texttt{x\_scale} & \texttt{jac} \\
\texttt{ftol} &  10{-8}\\
\texttt{xtol} &  10{-8}\\
\texttt{gtol} &  10{-8}\\
\texttt{f\_scale} &  0.001\\
\texttt{max\_nfev} &  10000\\
\texttt{loss} &  \texttt{soft\_l1}\\
\texttt{method} & \texttt{trf}
\end{tabular}
\caption{Parameters for the SciPy.optimize least\_squares algorithm in the second optimization round}
\label{parsopt2}
\end{table}

\subsubsection{Normalising residuals}\label{normalisation}

We obtain one scale value per frequency in either of the two types of data used. For SOFC data, since each dataset is comprised of
several measurements across 233 frequencies, we take a median of impedances. This leaves us with a 233 point dataset. To calculate the scales, we
take the median absolute deviation (MAD) of impedances of each frequency.

For the battery data, we first need to derive the impedances from the raw data available. Since there are many datapoints available we randomly
sample 300 points for the data

We noticed that in some cases the scales can be quite small, which makes dividing residual values (which can also be small) with them unstable.
To this end, we replace the smallest 10 \% of the values of the scale with a fixed value. It should be noted that we calculate separate scale values
for imaginary and real parts of the impedances.

%In practice

\section{Results}

We apply methods described above to data from both solid oxide fuel cells, as well as batteries. We try to distinguish which of the numerical adaptations
of the algorithm (additional settings for the SciPy's least-squares solver, normalising residuals) have the largest impact on the fit.

The plots presented in the graphs are the Nyquist graphs of the 5 circuits which we fitted until convergence, ordered by their BIC values.

\subsection{SOFC data}

We took \texttt{eis\_20200227\_070006.npz} dataset from repository \\\href{https://portal.ijs.si/nextcloud/s/xTa2cmtfxXn2jSz}{https://portal.ijs.si/nextcloud/s/xTa2cmtfxXn2jSz}.
As mentioned in Section \ref{normalisation}, we have several repeated measurements of impedance at each frequency. We take the median of them to get 
one impedance value per frequency.

In Figure \ref{sofc4el} we see results of fitting circuits with up to 4 elements. They do not capture any of the details of the data.

\begin{figure}[h]
\centering
\includegraphics[width=\textwidth]{../../Results/Plots/Nyquists/eis_20200227_070006_no_G_relaxed_4_best_5.png}
\caption{Nyquist graphs of 5 best fitting circuits according to our algorithm with up to 4 elements}
\label{sofc4el}
\end{figure}
%Nadaljeval sem z ocenjevanjem parametrov, vendar pa tudi s 5 elementi rezultati niso bili najboljši, kot je razvidno iz Grafa \ref{sofc5el_nonorm}.

Results with an additional circuit element are only marginally better (Figure \ref{sofc5el_nonorm}).

\begin{figure}[h!]
\centering
\includegraphics[width=\textwidth]{../../Results/Plots/Nyquists/eis_20200227_070006_no_G_relaxed_5_best_5_fixed_starting_values_no_normalization.png}
\caption{Nyquist graphs of 5 best fitting circuits according to our algorithm with up to 5 elements}
\label{sofc5el_nonorm}
\end{figure}

Before expanding the search space with even more circuits, we decided to try to make the algorithm more robust.
We first implemented normalising residuals at each step of the algorithm and present results in Figure \ref{sofc5el}. There is noticable
improvement in how closely the graphs fit the data. When inspecting the results in Table \ref{fig:first-table-circuits} more closely, we see that
several parameters are close to their limits (especially $\alpha$). 

\begin{figure}[h!]
\centering
\includegraphics[width=\textwidth]{../../Results/Plots/Nyquists/eis_20200227_070006_no_G_relaxed_5_best_5.png}
\caption{Nyquist graphs of 5 best fitting circuits according to our algorithm with up to 5 elements and normalised residuals.}
\label{sofc5el}
\end{figure}

To improve on unstable solutions we robustify the algorithm with adaptations described in Section \ref{numalg}. This way we get some more diversity
in results in Graph \ref{sofc5elnum}.

\begin{figure}[h]
\centering
\includegraphics[width=\textwidth]{../../Results/Plots/Nyquists/eis_20200227_070006_no_G_relaxed_5_best_5_extra_numerics.png}
\caption{Nyquist graphs of 5 best fitting circuits according to our algorithm with up to 5 elements and robustified algorithm}
\label{sofc5elnum}
\end{figure}

Comparing the circuits obtained with the standard and robustified numerical procedures we get
can see that the former contains one parameter less (but same number of elements),
and none of the parameters are on the border of their permitted range. MSE difference
between them was minimal. The circuits in question were:

\[
\begin{aligned}
\textnormal{Robustified approach}
    &= R_s
    \seriesconn
    \left[
        \CPEel{Q_1}{\alpha_1}
        \parallel
        \left(
            R_2
            \seriesconn
            \left[
                R_3
                \parallel
                \CPEel{Q_4}{\alpha_4}
            \right]
        \right)
    \right],
    \\[1.5ex]
\textnormal{Standard approach}
    &= R_s
    \seriesconn
    \left[
        R_1
        \parallel
        \left(
            \CPEel{Q_2}{\alpha_2}
            \seriesconn
            \left[
                \CPEel{Q_3}{\alpha_3}
                \parallel
                \CPEel{Q_4}{\alpha_4}
            \right]
        \right)
    \right].
\end{aligned}
\]

\begin{figure}[htbp]
    \centering

    %------------------------------------------------
    % Model A
    %------------------------------------------------
    \begin{subfigure}[t]{0.47\textwidth}
        \centering
        \resizebox{0.95\linewidth}{!}{%
        \begin{circuitikz}[line width=0.7pt]

            % Series resistance
            \draw
            (0,0) node[ocirc, label=left:$a$] {}
            to[R, l={$R_s$}, fill=gray!10]
            (1.6,0)
            -- (2,0) coordinate (A);

            % Outer upper branch: CPE1
            \draw
            (A)
            -- (2,1)
            to[generic,
               l^={$\mathrm{CPE}(Q_1,\alpha_1)$},
               fill=blue!12]
            (6.6,1)
            -- (6.6,0) coordinate (B);

            % Outer lower branch: R2
            \draw
            (A)
            -- (2,-1)
            to[R, l_={$R_2$}, fill=gray!10]
            (3.6,-1)
            -- (3.9,-1) coordinate (C);

            % Nested upper branch: R3
            \draw
            (C)
            to[R, l^={$R_3$}, fill=gray!10]
            (5.6,-1) coordinate (D);

            % Nested lower branch: CPE4
            \draw
            (C)
            -- (3.9,-2.2)
            to[generic,
               l_={$\mathrm{CPE}(Q_4,\alpha_4)$},
               fill=blue!12]
            (5.6,-2.2)
            -- (D);

            % Rejoin outer branch
            \draw
            (D)
            -- (6.6,-1)
            -- (B);

            % Output
            \draw
            (B)
            -- (7.2,0)
            node[ocirc, label=right:$b$] {};

            % Junctions
            \fill (A) circle (1.7pt);
            \fill (B) circle (1.7pt);
            \fill (C) circle (1.7pt);
            \fill (D) circle (1.7pt);

        \end{circuitikz}%
        }

        \caption{Result of the robustified approach}
        \label{fig:first-table-model-a}
    \end{subfigure}
    \hfill
    %------------------------------------------------
    % Model B
    %------------------------------------------------
    \begin{subfigure}[t]{0.47\textwidth}
        \centering
        \resizebox{0.95\linewidth}{!}{%
        \begin{circuitikz}[line width=0.7pt]

            % Series resistance
            \draw
            (0,0) node[ocirc, label=left:$a$] {}
            to[R, l={$R_s$}, fill=gray!10]
            (1.6,0)
            -- (2,0) coordinate (A);

            % Outer upper branch: R1
            \draw
            (A)
            -- (2,1)
            to[R, l^={$R_1$}, fill=gray!10]
            (6.6,1)
            -- (6.6,0) coordinate (B);

            % Outer lower branch: CPE2
            \draw
            (A)
            -- (2,-1)
            to[generic,
               l_={$\mathrm{CPE}(Q_2,\alpha_2)$},
               fill=blue!12]
            (3.6,-1)
            -- (3.9,-1) coordinate (C);

            % Nested upper branch: CPE3
            \draw
            (C)
            to[generic,
               l^={$\mathrm{CPE}(Q_3,\alpha_3)$},
               fill=blue!12]
            (5.6,-1) coordinate (D);

            % Nested lower branch: CPE4
            \draw
            (C)
            -- (3.9,-2.2)
            to[generic,
               l_={$\mathrm{CPE}(Q_4,\alpha_4)$},
               fill=blue!12]
            (5.6,-2.2)
            -- (D);

            % Rejoin outer branch
            \draw
            (D)
            -- (6.6,-1)
            -- (B);

            % Output
            \draw
            (B)
            -- (7.2,0)
            node[ocirc, label=right:$b$] {};

            % Junctions
            \fill (A) circle (1.7pt);
            \fill (B) circle (1.7pt);
            \fill (C) circle (1.7pt);
            \fill (D) circle (1.7pt);

        \end{circuitikz}%
        }

        \caption{Result of the standard approach}
        \label{fig:first-table-model-b}
    \end{subfigure}

    \caption{Resulting circuits from two approaches with parameters described in \ref{tab:circuit_comparison_sofc}:
    (a) robustified algorithm;
    (b) standard algorithm.}
    \label{fig:first-table-circuits}
\end{figure}


\begin{table}[htbp]
\centering
\caption{Parameter comparison of the two circuits}
\label{tab:circuit_comparison_sofc}
\small
\begin{tabularx}{\textwidth}{@{}lXX@{}}
\toprule
\textbf{Parameter} & \textbf{Robustified algorithm} & \textbf{Standard approach} \\
\midrule

%Circuit &
%$\displaystyle
%R_s +
%\left[
%\mathrm{CPE}(Q_1,\alpha_1)
%\parallel
%\left(
%R_2 +
%\left[
%R_3 \parallel \mathrm{CPE}(Q_4,\alpha_4)
%\right]
%\right)
%\right]$
%&
%$\displaystyle
%R_s +
%\left[
%R_1
%\parallel
%\left(
%\mathrm{CPE}(Q_2,\alpha_2) +
%\left[
%\mathrm{CPE}(Q_3,\alpha_3)
%\parallel
%\mathrm{CPE}(Q_4,\alpha_4)
%\right]
%\right)
%\right]$
%\\[2ex]

Element count     & 5 & 5 \\
Parameter count   & 7 & 8 \\
%Random start      & Yes & Yes \\
MSE               & $1.89647\times10^{-6}$ & $1.64223\times10^{-6}$ \\
BIC & $-3031.73789$ & $-3059.82487$ \\
%$\mathrm{lp\_error}$ & 0.865611 & Not reported \\

\midrule
\multicolumn{3}{@{}l}{\textbf{Circuit parameters}} \\
\midrule

$R_s$ (\si{\ohm}) &
0.0463228 &
0.0461368 \\

$R_1$ (\si{\ohm}) &
--- &
0.0258082 \\

$R_2$ (\si{\ohm}) &
0.0129929 &
--- \\

$R_3$ (\si{\ohm}) &
0.0133333 &
--- \\

$Q_1$ (\si{\siemens\second\tothe{\alpha_1}}) &
0.542887 &
--- \\

$\alpha_1$ &
0.795783 &
--- \\

$Q_2$ (\si{\siemens\second\tothe{\alpha_2}}) &
--- &
1.05020 \\

$\alpha_2$ &
--- &
0.990000 \\

$Q_3$ (\si{\siemens\second\tothe{\alpha_3}}) &
--- &
50.8970 \\

$\alpha_3$ &
--- &
0.0298022 \\

$Q_4$ (\si{\siemens\second\tothe{\alpha_4}}) &
3.33826 &
0.469656 \\

$\alpha_4$ &
0.941526 &
0.819800 \\

\midrule
%\multicolumn{3}{@{}l}{\textbf{Rezistenčne lastnosti}} \\
%\midrule
%
%Teoretičen upor vezja pri visokih frekvencah, $R_{\infty}$ (\si{\ohm}) &
%0.0463228 &
%0.0461368 \\
%
%Teoretičen upor vezja pri nizkih frekvencah, $R_0$ (\si{\ohm}) &
%0.0726489 &
%0.0719450 \\
%
%Razlika med obema, $R_{\mathrm{pol}}$ (\si{\ohm}) &
%0.0263261 &
%0.0258082 \\
%
%\bottomrule
\end{tabularx}
\end{table}






\clearpage



\subsection{Battery data}

We apply the methods to data from batteries as well. We used data from 
\href{https://portal.ijs.si/nextcloud/s/2TRZjgGGLZygtNs}{https://portal.ijs.si/nextcloud/s/2TRZjgGGLZygtNs}.

%Enake metode sem uporabil tudi na podatkih iz baterij iz repozitorija \\\href{https://portal.ijs.si/nextcloud/s/2TRZjgGGLZygtNs}{https://portal.ijs.si/nextcloud/s/2TRZjgGGLZygtNs}.
%Tej so predvsem v nizkih frekvencah bolj razpršeni, zato je grafe vizualno morda težje ocenjevati.



\begin{figure}[h]
\centering
\includegraphics[width=\textwidth]{../../Results/Plots/Nyquists/bank3_20260208-084547_1_no_G_relaxed_4_best_5.png}
\caption{Nyquist graphs of 5 best fitting circuits according to our algorithm with up to 4 elements}
\end{figure}



\begin{figure}[h]
\centering
\includegraphics[width=\textwidth]{../../Results/Plots/Nyquists/bank3_20260208-084547_1_no_G_relaxed_5_best_5.png}
\caption{Nyquist graphs of 5 best fitting circuits according to our algorithm with up to 5 elements}\end{figure}

\begin{figure}[h]
\centering
\includegraphics[width=\textwidth]{../../Results/Plots/Nyquists/bank3_20260208-084547_1_no_G_relaxed_5_best_5_extra_numerics.png}
\caption{Nyquist graphs of 5 best fitting circuits according to our algorithm with up to 5 elements and robustified algorithm}
\end{figure}

We compare the 5 element circuits obtained from robustified and standard numerical algorithm. 

\[
\begin{aligned}
\textnormal{Robustified approach}
    &= R_s
    \seriesconn
    \left[
        \left(
            \CPEel{Q_1}{\alpha_1}
            \seriesconn
            \CPEel{Q_2}{\alpha_2}
        \right)
        \parallel
        \left(
            R_3
            \seriesconn
            \CPEel{Q_4}{\alpha_4}
        \right)
    \right],
    \\[1.5ex]
\textnormal{Standard approach}
    &= R_s
    \seriesconn
    \CPEel{Q_1}{\alpha_1}
    \seriesconn
    \left[
        R_2
        \parallel
        \left(
            \CPEel{Q_3}{\alpha_3}
            \seriesconn
            \CPEel{Q_4}{\alpha_4}
        \right)
    \right].
\end{aligned}
\]

\begin{figure}[htbp]
    \centering

    %------------------------------------------------
    % Extra-numerics circuit
    %------------------------------------------------
    \begin{subfigure}[t]{0.47\textwidth}
        \centering
        \resizebox{0.95\linewidth}{!}{%
        \begin{circuitikz}[line width=0.7pt]

            % Series resistance
            \draw
            (0,0) node[ocirc, label=left:$a$] {}
            to[R, l={$R_s$}, fill=gray!10]
            (1.6,0)
            -- (2,0) coordinate (A);

            % Upper branch: CPE1 + CPE2
            \draw
            (A)
            -- (2,1)
            to[generic,
               l^={$\mathrm{CPE}(Q_1,\alpha_1)$},
               fill=blue!12]
            (4.2,1)
            to[generic,
               l^={$\mathrm{CPE}(Q_2,\alpha_2)$},
               fill=blue!12]
            (6.4,1)
            -- (6.4,0) coordinate (B);

            % Lower branch: R3 + CPE4
            \draw
            (A)
            -- (2,-1)
            to[R, l_={$R_3$}, fill=gray!10]
            (4.2,-1)
            to[generic,
               l_={$\mathrm{CPE}(Q_4,\alpha_4)$},
               fill=blue!12]
            (6.4,-1)
            -- (B);

            % Output
            \draw
            (B)
            -- (7,0)
            node[ocirc, label=right:$b$] {};

            % Junctions
            \fill (A) circle (1.7pt);
            \fill (B) circle (1.7pt);

        \end{circuitikz}%
        }

        \caption{Extra-numerics fit}
        \label{fig:second-table-extra}
    \end{subfigure}
    \hfill
    %------------------------------------------------
    % Standard circuit
    %------------------------------------------------
    \begin{subfigure}[t]{0.47\textwidth}
        \centering
        \resizebox{0.95\linewidth}{!}{%
        \begin{circuitikz}[line width=0.7pt]

            % Rs and CPE1 in series
            \draw
            (0,0) node[ocirc, label=left:$a$] {}
            to[R, l={$R_s$}, fill=gray!10]
            (1.4,0)
            to[generic,
               l^={$\mathrm{CPE}(Q_1,\alpha_1)$},
               fill=blue!12]
            (3.7,0)
            -- (4,0) coordinate (A);

            % Upper branch: R2
            \draw
            (A)
            -- (4,1)
            to[R, l^={$R_2$}, fill=gray!10]
            (7.1,1)
            -- (7.1,0) coordinate (B);

            % Lower branch: CPE3 + CPE4
            \draw
            (A)
            -- (4,-1)
            to[generic,
               l_={$\mathrm{CPE}(Q_3,\alpha_3)$},
               fill=blue!12]
            (5.5,-1)
            to[generic,
               l_={$\mathrm{CPE}(Q_4,\alpha_4)$},
               fill=blue!12]
            (7.1,-1)
            -- (B);

            % Output
            \draw
            (B)
            -- (7.7,0)
            node[ocirc, label=right:$b$] {};

            % Junctions
            \fill (A) circle (1.7pt);
            \fill (B) circle (1.7pt);

        \end{circuitikz}%
        }

        \caption{Standard fit}
        \label{fig:second-table-standard}
    \end{subfigure}

    \caption{Side by side cirucit comparison:
    (a) Robustified approach
    (b) Standard approach}
    \label{fig:second-table-circuits}
\end{figure}

\begin{table}[htbp]
\centering
\caption{5 element circuit parameter comparison}
\label{tab:circuit_comparison}
\small
\begin{tabularx}{\textwidth}{@{}lXX@{}}
\toprule
\textbf{Parameter}
& \textbf{Robustified approach}
& \textbf{Standard approach} \\
\midrule

%Circuit
%&
%$\begin{aligned}
%R_s + \bigl[
%&\bigl(
%\operatorname{CPE}(Q_1,\alpha_1)
%+\operatorname{CPE}(Q_2,\alpha_2)
%\bigr)\\
%&\parallel
%\bigl(
%R_3+\operatorname{CPE}(Q_4,\alpha_4)
%\bigr)
%\bigr]
%\end{aligned}$
%&
%$\begin{aligned}
%R_s &+ \operatorname{CPE}(Q_1,\alpha_1)\\
%&+ \bigl[
%R_2 \parallel
%\bigl(
%\operatorname{CPE}(Q_3,\alpha_3)
%+\operatorname{CPE}(Q_4,\alpha_4)
%\bigr)
%\bigr]
%\end{aligned}$
%\\[2ex]

Element count
& 5
& 5 \\

Parameter count
& 8
& 8 \\

MSE
& $2.72147\times10^{-8}$
& $2.33564\times10^{-8}$ \\

BIC
& $-5180.22$
& $-5226.09$ \\

\midrule
\multicolumn{3}{@{}l}{\textbf{Circuit parameters}} \\
\midrule

%RMSE
%& $1.64969\times10^{-4}$
%& $1.52828\times10^{-4}$ \\
%
%
%Function evaluations
%& 100
%& 62 \\
%
%Termination
%& Maximum number of evaluations exceeded
%& \texttt{ftol} condition satisfied \\


$R_s$ (\(\Omega\))
& 0.0160873
& $8.08927\times10^{-5}$ \\

$R_2$ (\(\Omega\))
& ---
& 0.0225741 \\

$R_3$ (\(\Omega\))
& 0.00747417
& --- \\

$Q_1$ (\(\mathrm{S\,s}^{\alpha_1}\))
& 122.545
& 505.257 \\

$\alpha_1$
& 0.664787
& 0.612086 \\

$Q_2$ (\(\mathrm{S\,s}^{\alpha_2}\))
& 5.36106
& --- \\

$\alpha_2$
& 0.640431
& --- \\

$Q_3$ (\(\mathrm{S\,s}^{\alpha_3}\))
& ---
& 4.29185 \\

$\alpha_3$
& ---
& 0.166491 \\

$Q_4$ (\(\mathrm{S\,s}^{\alpha_4}\))
& 859.752
& 0.122094 \\

$\alpha_4$
& 0.693547
& 0.989999 \\

\bottomrule
\end{tabularx}
\end{table}

%\begin{figure}
%\centering
%\begin{subfigure}{.5\textwidth}
%  \centering
%  \includegraphics[width=.4\linewidth]{../../Results/Plots/Nyquists/eis_20200227_070006_no_G_relaxed_4_best_5.png}
%  \caption{A subfigure}
%  \label{fig:sub1}
%\end{subfigure}%
%\begin{subfigure}{.5\textwidth}
%  \centering
%  \includegraphics[width=.4\linewidth]{../../Results/Plots/Nyquists/eis_20200227_070006_no_G_relaxed_5_best_5.png}
%  \caption{A subfigure}
%  \label{fig:sub2}
%\end{subfigure}
%\caption{A figure with two subfigures}
%\label{fig:test}
%\end{figure}
\clearpage
\nocite{*} 
\printbibliography
%\bibliographystyle{plainnat}

%\bibliography{references} 


\end{document}
